\documentclass[10pt,a4paper]{amsart}
\usepackage[margin=19mm]{geometry}
\usepackage{amsmath,amssymb,mathtools}
\usepackage[scr=rsfs,cal=euler]{mathalfa}
\usepackage{xfrac}
\usepackage[unicode,hidelinks,bookmarks=false]{hyperref}
\newcommand{\ie}{i.e.\ }
\newcommand{\cf}{cf.\ }
\newcommand{\resp}{respectively\ }
\newcommand{\Subsection}{\subsection}
\providecommand{\nobreakdash}{\nobreak\hbox{-}}
\setlength{\emergencystretch}{2em}
\multlinegap0pt
\usepackage{amssymb}
\usepackage{xcolor}
\usepackage{tikz}
\newtheorem{lemma}{Lemma}
\newtheorem{theorem}{Theorem}
\newcommand{\Adem}{A_{\mathrm{dem}}}
\title{Optimal pebbling of the hypercube}
\author{Lior Pachter}
\date{Source: arXiv:2606.01685v2 (CC BY 4.0)}
\begin{document}
\maketitle

\noindent\emph{Source and licence.} Optimal pebbling of the hypercube. Version arXiv:2606.01685v2 (CC BY 4.0), distributed by arXiv under the Creative Commons Attribution 4.0 International licence, http://creativecommons.org/licenses/by/4.0/, as recorded in the arXiv licence metadata for this exact version and shown on the abstract page https://arxiv.org/abs/2606.01685v2. Full licence notice: \texttt{licenses/arxiv-2606.01685-CC-BY-4.0.txt}.\par\smallskip

\noindent\textit{Editorial scope.} Counted unit: the article's theorem thm:upper (source0.tex lines 332--334). The article's complete original proof of it is delivered with it (source0.tex lines 335--396, one \texttt{proof} environment of 62 lines). No mathematics is written, repaired or re-derived: the statement and the proof are the article's own lines, in the article's own notation, and no retained line is substituted or re-flowed.  Retained uncounted as the source-local context the proof needs: lines 143--155; lines 252--258.  Nothing was omitted from any retained span, so the package carries no omission record.  No retained line contains a citation, so the wrapper carries no bibliography and no bibliography entry.  No retained line contains a figure environment, an image-inclusion command or drawing code, so no image is delivered and the package declares no asset dependency.  Editorial formatting only, in addition: an AMS \texttt{amsart} preamble that restates the article's own declarations and macros used by the retained lines, namely the environments \texttt{lemma}, \texttt{theorem} on separate counters, together with the standard packages those lines need.  The retained environments print in this document as Lemma 1, Theorem 1 in order of appearance, whereas the article prints the same items with the numbering of its own full document.  The complete native source of the article as distributed in the arXiv e-print for this exact version is bundled unchanged as \texttt{sources/201846/source0.tex}.
\begin{lemma}[Near-optimal high-demand lemma]\label{lem:high-demand}
There are constants $\Adem,C,n_0>0$ such that the following holds.  If $n\ge n_0$,
$n^{-3}\le \delta\le 1/2$, and
\[
T\,\ge\, C\delta^{-2}n2^{\Adem\sqrt{n\log(2/\delta)}},
\]
then
\[
o_T(Q_n)\,\le\, (1+\delta)T\left(\frac43\right)^n.
\]
Moreover, the witnessing distribution may be chosen with every occupied pile of size at
least $2^{\lfloor n/3\rfloor}$.
\end{lemma}

\begin{lemma}[Product / composition step]\label{lem:product-step}
Suppose $E$ is a solvable distribution on $Q_m$, and for each occupied vertex $z$ of $Q_m$
with $E(z)=s_z$ one has $o_{s_z}(Q_a)\le (1+\delta)s_z\left(\frac43\right)^a$.  Then
\[
o(Q_n)\,\le\, (1+\delta)|E|\left(\frac43\right)^a.
\]
\end{lemma}

\begin{theorem}[Upper bound]\label{thm:upper}
$o(Q_n)\,=\,O\!\left(\left(\frac43\right)^n\right).$
\end{theorem}

\begin{proof}
Let $\Adem,C$ be the constants from Lemma~\ref{lem:high-demand}. In the recursion, the
smaller factor $Q_m$ supplies the demand values $s_z$, while the larger factor $Q_a$
supplies those demands inside the fibers.  The inductive pile-size lower bound will give
$s_z\ge 2^{\lfloor m/5\rfloor}$, so $m$ must be large enough for
$2^{\lfloor m/5\rfloor}$ to exceed the high-demand threshold in dimension $a$.
Choose $K>5\Adem\sqrt 2$.  Then $m/5$ dominates
$\Adem\sqrt{n\log(2n^2)}$ for large $n$; the floor only changes constants.
Choose $N_0$ so large
that, for every $n\ge N_0$, with
\[
m\,=\,\left\lceil K\sqrt{n\log n}\right\rceil,\qquad a\,=\,n-m,
\]
we have $m\le n/4$,
\[
2^{\lfloor m/5\rfloor}\,\ge\, C n^5 2^{\Adem\sqrt{n\log(2n^2)}},
\]
and $a\ge n_0$ and $a^{-3}\le n^{-2}\le 1/2$, so the high-demand lemma applies in
dimension $a$ with $\delta=n^{-2}$.
\label{def:parameters}

We define solvable distributions $D_n$ recursively.  For every $n<N_0$, choose any
solvable distribution $D_n$ on $Q_n$ whose occupied piles have size at least
$2^{\lfloor n/5\rfloor}$; for example, place $2^{\lfloor n/5\rfloor}$ pebbles on every
vertex.  For $n\ge N_0$, assume
that $D_m$ has already been constructed for
$m=\left\lceil K\sqrt{n\log n}\right\rceil$.
For each occupied vertex $z$ of $D_m$, let $s_z=D_m(z)$.  By induction every occupied
pile of $D_m$ satisfies $s_z\ge 2^{\lfloor m/5\rfloor}$, and by the choice of
$K,N_0$ this is at least $C\delta^{-2}a2^{\Adem\sqrt{a\log(2/\delta)}}$.
Thus Lemma~\ref{lem:high-demand} applies on the $Q_a$ factor for every required demand
$s_z$.  Define $D_n$ on $Q_n=Q_a\square Q_m$ by putting, in the fiber $Q_a\times\{z\}$, an
$s_z$-demand distribution on $Q_a$ of size at most $(1+n^{-2})s_z\left(\frac43\right)^a$.
By Lemma~\ref{lem:product-step} the distribution $D_n$ is solvable, and
\[
|D_n|\,\le\, (1+n^{-2})|D_m|\left(\frac43\right)^a.
\]
The pile-size invariant is preserved: the high-demand lemma uses piles of size at least
$2^{\lfloor a/3\rfloor}$, and since $a\ge 3n/4$ we have
$\lfloor a/3\rfloor\ge n/4-1\ge n/5\ge \lfloor n/5\rfloor$ whenever $n\ge 20$.
\label{lem:recurrence}

Put $R_n=|D_n|/(4/3)^n$.  The construction yields $R_n\le (1+n^{-2})R_m$, and for the
fully integral recursive bound we use the harmless weaker recurrence
$R_n\le (1+2n^{-2})R_m$, which absorbs the ceiling in the integer fiber multiplier and
changes only the final absolute constant.  Iterate this recurrence while $n_i\ge N_0$,
where $n_{i+1}=\left\lceil K\sqrt{n_i\log n_i}\right\rceil$.
By increasing $N_0$ we may assume $n_{i+1}\le n_i/4$ throughout the iteration, so that
$n_{i+1}^{-2}\ge 16\,n_i^{-2}$ and the sequence $(n_i^{-2})$ grows at least geometrically.
Summing this geometric series, $\sum_i n_i^{-2}\le n_{i_*}^{-2}\cdot\frac{16}{15}\le
\frac{16}{15N_0^2}$, where $n_{i_*}\ge N_0$ is the last active term in the iteration.
At the terminal dimension $b<N_0$, the base construction gives
$R_b\le 2^{\lfloor b/5\rfloor}\left(\frac32\right)^b$, so all terminal costs are bounded by
$R_{\mathrm{base}}=\sum_{j=0}^{N_0-1}2^{\lfloor j/5\rfloor}\left(\frac32\right)^j$.  Hence
\[
R_n\,\le\, R_{\mathrm{base}}\prod_i(1+2n_i^{-2})
\,\le\, R_{\mathrm{base}}\exp\!\left(2\sum_i n_i^{-2}\right)
\,\le\, R_{\mathrm{base}}\exp\!\left(\frac{4}{N_0^2}\right).
\]
\label{lem:loss}
This is an absolute constant, and hence $o(Q_n)\le |D_n|=O\!\left(\left(\frac43\right)^n\right)$.
\end{proof}

\end{document}
